% =========================================================================
% SUMÁRIO (ANTES DO RESUMO)
% =========================================================================
\pagestyle{empty}
\renewcommand{\contentsname}{Sumário}
\setcounter{tocdepth}{1}
\tableofcontents
\clearpage

% =========================================================================
% RESUMO E INTRODUÇÃO (COMPARTILHAM UMA ÚNICA PÁGINA)
% =========================================================================
\pagestyle{textual}
% =========================================================================
% SEÇÃO 1: RESUMO
% =========================================================================
\section{Resumo}\label{sec:resumo}

Fluxos em digrafos constituem uma área fundamental da Otimização Combinatória, com amplas aplicações práticas em redes de transporte e alocação de recursos. Este projeto adota uma abordagem que combina fundamentos teóricos, implementação algorítmica e validação experimental para o estudo de problemas de fluxo máximo e fluxo de custo mínimo (veja~\cite{ahuja1993network}). Foram estudados e implementados algoritmos clássicos e modernos de fluxo máximo e de fluxo de custo mínimo, bem como técnicas de modelagem e reduções combinatórias, incluindo a aplicação prática em Segmentação Interativa de Imagens via corte mínimo. A validação prática foi realizada através de experimentos computacionais utilizando instâncias de \textit{benchmarks} reconhecidas da literatura (veja~\cite{Jensen2022}).


\vspace{0.35cm}

% =========================================================================
% SEÇÃO 2: INTRODUÇÃO
% =========================================================================
\section{Introdução}\label{sec:introducao}

A Otimização Combinatória é uma área de pesquisa situada na interseção entre Matemática e Ciência da Computação (veja~\cite{korte2018combinatorial}). Do ponto de vista matemático, a área mantém forte influência com a Teoria dos Grafos, uma vez que digrafos constituem objetos combinatórios com amplo poder de modelagem (veja~\cite{bondy2008graph, diestel2017graph}).
Sob a perspectiva computacional, o foco reside na obtenção de algoritmos eficientes que produzam soluções garantidamente ótimas (veja~\cite{schrijver2003combinatorial, cook1998combinatorial}). Dentre os tópicos de destaque, os problemas de fluxos em digrafos ocupam posição central, tratando de redes onde recursos devem ser transportados respeitando restrições de capacidade (veja~\cite{ahuja1993network}).
O objetivo deste projeto é proporcionar uma formação sólida na área, integrando teoria e prática, abrangendo desde fundamentos conceituais até implementações de alta performance em C++23 e aplicações reais em Visão Computacional.

O problema do $st$-fluxo máximo, onde se busca maximizar a quantidade de fluxo enviada de um vértice fonte $s$ para um vértice de destino $t$, é um dos pilares deste estudo (veja~\cite{ford1956maximal}). Foram analisados e implementados algoritmos clássicos e eficientes, como Ford-Fulkerson, Edmonds-Karp, Dinic e Push-Relabel com Heurística de Gap (veja~\cite{ford1956maximal, edmonds1972theoretical, dinic1970algorithm, goldberg1988new}).
O projeto contempla também o problema de fluxo de custo mínimo, que busca minimizar o custo total enquanto satisfaz as demandas da rede, com destaque para os algoritmos de cancelamento de ciclos de Klein, caminhos sucessivos e o \textit{network simplex}, uma especialização do algoritmo simplex para redes (veja~\cite{dantzig1951application, ahuja1993network}).
A fundamentação teórica desenvolvida reafirma que muitos problemas combinatórios, como o emparelhamento máximo em digrafos bipartidos, coberturas mínimas e segmentação de imagens via corte mínimo, podem ser resolvidos via reduções para problemas de fluxo (veja~\cite{ahuja1993network, boykov2001interactive}).

\clearpage

% =========================================================================
% METODOLOGIA (MÁXIMO 1 PÁGINA)
% =========================================================================
% =========================================================================
% SEÇÃO 2: METODOLOGIA
% =========================================================================
\section{Metodologia}\label{sec:metodologia}

\subsection{Materiais e Métodos}

A metodologia da pesquisa desenvolveu-se em ciclos de fundamentação teórica, modelagem matemática, implementação computacional e experimentação empírica sistemática:

\begin{enumerate}[leftmargin=1.5em, itemsep=0.15em, topsep=0.2em, parsep=0em]
    \item \textbf{Ambiente de Software e Compilação:} As implementações foram desenvolvidas em linguagem C++23, compiladas com GCC 13+ sob nível de otimização \lstinline{-O2}. Para garantir segurança contra estouros de pilha (\textit{stack overflow}) em buscas em profundidade recursivas sobre grafos densos, fixou-se a pilha do executável em 256~MB através da flag \lstinline{-Wl,-z,stack-size=268435456};
    \item \textbf{Estruturas de Dados e Localidade de Cache:} As redes de fluxo foram estruturadas em listas de adjacência indexadas por vetores contíguos em memória (\lstinline{std::vector}), eliminando alocações dinâmicas dispersas de nós e ponteiros e favorecendo a taxa de acertos na memória cache da CPU;
    \item \textbf{Protocolo Experimental e Isolamento de CPU:} Os testes foram cronometrados exclusivamente durante o cálculo algorítmico por meio do relógio monotônico \lstinline{std::chrono::high_resolution_clock}, descartando o tempo de leitura de disco e parsing das instâncias;
    \item \textbf{Bases de Teste Padronizadas (DIMACS):} Utilizaram-se as famílias de teste do \textit{First DIMACS Challenge}~\cite{Jensen2022, ahuja1993network}, geradas oficialmente pelos programas de referência de Rutgers (famílias Washington e Genrmf para fluxo máximo; família Netgen para fluxo de custo mínimo).
\end{enumerate}

\subsection{Etapas da Pesquisa e Cronograma de Atividades}

As atividades planejadas no plano de trabalho da Iniciação Científica (vigência 2025/2026) foram estruturadas conforme o cronograma oficial:

\begin{itemize}[leftmargin=1.5em, itemsep=0.15em, topsep=0.2em, parsep=0em]
    \item \textbf{Atividade 1:} Estudo dos fundamentos de Otimização Combinatória e Teoria dos Grafos~\cite{cook1998combinatorial}. \textbf{[Concluída]}
    \item \textbf{Atividade 2:} Estudo teórico do problema de fluxo máximo e algoritmos clássicos~\cite{ford1956maximal, edmonds1972theoretical, dinic1970algorithm, goldberg1988new}. \textbf{[Concluída]}
    \item \textbf{Atividade 3:} Implementação em C++ dos algoritmos de fluxo máximo (Ford-Fulkerson, Edmonds-Karp, Dinic, Push-Relabel e Push-Relabel com Gap Heuristic). \textbf{[Concluída]}
    \item \textbf{Atividade 4:} Estudo de técnicas de modelagem e reduções de problemas combinatórios. \textbf{[Concluída]}
    \item \textbf{Atividade 5:} Implementação de problemas modelados (Emparelhamento Bipartido Máximo, Cobertura Mínima, Conjunto Independente e Caminhos Disjuntos). \textbf{[Concluída]}
    \item \textbf{Atividade 6:} Elaboração do relatório de Iniciação Científica (este documento). \textbf{[Concluída]}
    \item \textbf{Atividade 7:} Estudo teórico do problema de fluxo de custo mínimo, dualidade e condições de otimalidade por ciclos negativos e custos reduzidos~\cite{ahuja1993network}. \textbf{[Concluída]}
    \item \textbf{Atividade 8:} Implementação em C++ dos algoritmos de custo mínimo (Cycle Canceling, Successive Shortest Path com SPFA e Dijkstra acelerado por potenciais, e \textit{Network Simplex}). \textbf{[Concluída]}
    \item \textbf{Atividade 9:} Baterias de testes experimentais em instâncias padronizadas de larga escala (DIMACS). \textbf{[Concluída]}
    \item \textbf{Atividade 10:} Análise comparativa de desempenho assintótico e empírico dos algoritmos implementados. \textbf{[Concluída]}
    \item \textbf{Atividade 11:} Aplicação prática em Segmentação de Imagens e consolidação dos resultados da pesquisa. \textbf{[Concluída]}
\end{itemize}

A Figura~\ref{fig:pipeline_metodologico} ilustra o fluxo de trabalho articulado ao longo da pesquisa.

\begin{figure}[!ht]
\centering
\resizebox{0.95\textwidth}{!}{%
\begin{tikzpicture}[thick, node distance=0.6cm]
  \node[stage box] (s1) at (0,0) {\textbf{1. Teoria \& Dualidade}\\Max-Flow Min-Cut\\Reduções Bipartidas\\Fluxo Custo Mínimo};
  \node[stage box] (s2) at (3.3,0) {\textbf{2. Arquitetura C++23}\\Herança \& Polimorfismo\\\texttt{FlowNetwork}\\\texttt{CostNetwork}};
  \node[stage box] (s3) at (6.6,0) {\textbf{3. Motores Algorítmicos}\\Edmonds-Karp, Dinic\\Push-Relabel (Heur.)\\SSP, Network Simplex};
  \node[stage box] (s4) at (9.9,0) {\textbf{4. Certificação Cruzada}\\Consistência de Fluxo\\Dualidade Forte\\Google Test};
  \node[stage box] (s5) at (13.2,0) {\textbf{5. Benchmarks}\\Grades 2D, Camadas\\Grafos Aleatórios\\Segmentação Imagens};

  \draw[->, line width=1.5pt, blue!80!black] (s1) -- (s2);
  \draw[->, line width=1.5pt, blue!80!black] (s2) -- (s3);
  \draw[->, line width=1.5pt, blue!80!black] (s3) -- (s4);
  \draw[->, line width=1.5pt, blue!80!black] (s4) -- (s5);
\end{tikzpicture}%
}
\caption{Fluxo metodológico integrando fundamentação teórica, engenharia de software em C++23, validação cruzada e testes empíricos.}
\label{fig:pipeline_metodologico}
\end{figure}

\clearpage

% =========================================================================
% FUNDAMENTAÇÃO, RESULTADOS E APLICAÇÃO
% =========================================================================
% =========================================================================
% SEÇÃO 3: FUNDAMENTAÇÃO TEÓRICA
% =========================================================================
\section{Fundamentação Teórica}\label{sec:fundamentacao}

Esta seção formaliza os modelos matemáticos, teoremas estruturais e mecanismos algorítmicos que regem os problemas de Fluxo Máximo e Fluxo de Custo Mínimo.

\subsection{Base Matemática dos Grafos de Fluxo}

Uma rede de fluxo é representada por um digrafo $D = (V, A)$, onde $V$ denota o conjunto de vértices e $A$ o conjunto de arcos direcionados (veja~\cite{ahuja1993network, bondy2008graph}). A rede possui dois vértices terminais distinguidos: uma fonte $s \in V$ e um sumidouro $t \in V$. Cada arco $(u, v) \in A$ possui uma capacidade não negativa $c(u, v) \ge 0$, estipulada pela função $c: A \to \mathbb{R}_{\ge 0}$, que estabelece o limite superior de tráfego admissível no arco.

Um fluxo na rede $D$ é uma função $f: A \to \mathbb{R}_{\ge 0}$ que associa a cada arco $(u, v)$ um escalar $f(u, v)$. O fluxo $f$ é dito viável se satisfaz duas condições axiomáticas (veja~\cite{ford1956maximal, cook1998combinatorial}):
\begin{enumerate}
    \item \textbf{Restrição de Capacidade:} O fluxo em cada conexão é não negativo e não excede sua capacidade:
    \begin{equation}
        0 \le f(u, v) \le c(u, v), \quad \forall (u, v) \in A
    \end{equation}
    \item \textbf{Conservação de Fluxo:} Para todo vértice intermediário $v \in V \setminus \{s, t\}$, a soma dos fluxos que entram em $v$ é igual à soma dos fluxos que dele saem:
    \begin{equation}
        \sum_{\{u \mid (u, v) \in A\}} f(u, v) = \sum_{\{w \mid (v, w) \in A\}} f(v, w)
    \end{equation}
\end{enumerate}

O valor total do fluxo, denotado por $|f|$, quantifica a taxa líquida que emerge da fonte $s$:
\begin{equation}
    |f| = \sum_{\{v \mid (s, v) \in A\}} f(s, v) - \sum_{\{u \mid (u, s) \in A\}} f(u, s)
\end{equation}
Pela conservação nodal, esse valor coincide rigorosamente com o fluxo líquido absorvido pelo sumidouro $t$. O Problema do Fluxo Máximo consiste em determinar uma função viável $f$ que maximize $|f|$.

\subsection{Redes Residuais, Caminhos Aumentantes e Dualidade}

Dado um fluxo viável $f$, a rede residual $D_f = (V, A_f)$ quantifica a capacidade ociosa dos arcos originais e a possibilidade de cancelamento de fluxos previamente despachados (veja~\cite{ahuja1993network}). A capacidade residual $c_f(u, v)$ é definida por:
\begin{equation}
    c_f(u, v) = c(u, v) - f(u, v) + f(v, u)
\end{equation}
O conjunto $A_f$ é formado por todos os arcos direcionados $(u, v)$ tais que $c_f(u, v) > 0$. Um $st$-caminho aumentante é um caminho direcionado simples de $s$ para $t$ em $D_f$. Se existir tal caminho $P$, a capacidade residual gargalo $\delta = \min_{(u, v) \in P} c_f(u, v) > 0$ pode ser enviada ao longo de $P$, gerando um novo fluxo viável $f'$ com valor $|f'| = |f| + \delta$.

Um $st$-corte $(S, T)$ é uma partição de $V$ tal que $s \in S$ e $t \in T = V \setminus S$. A capacidade do corte é a soma das capacidades dos arcos direcionados de $S$ para $T$:
\begin{equation}
    c(S, T) = \sum_{u \in S, v \in T, (u, v) \in A} c(u, v)
\end{equation}

\begin{prop}[Teorema do Fluxo Máximo e Corte Mínimo --- Ford-Fulkerson, 1956]\label{prop:mfmc_rel}
Em qualquer rede de fluxo, o valor máximo de um $st$-fluxo viável é igual à capacidade mínima de um $st$-corte:
\begin{equation}
    \max_{f \text{ viável}} |f| = \min_{(S, T) \text{ corte}} c(S, T)
\end{equation}
Além disso, as seguintes condições são equivalentes: (i) $f$ é um fluxo máximo; (ii) $D_f$ não contém nenhum $st$-caminho aumentante; (iii) existe um corte $(S^*, T^*)$ tal que $|f| = c(S^*, T^*)$.
\end{prop}

\noindent \textbf{Demonstração Construtiva:} Para qualquer fluxo viável $f$ e qualquer corte $(S, T)$, tem-se $|f| \le c(S, T)$. Se $D_f$ não possui caminho aumentante de $s$ a $t$, seja $S^*$ o conjunto de vértices alcançáveis a partir de $s$ em $D_f$ e $T^* = V \setminus S^*$. Como não há caminho de $s$ a $t$, $s \in S^*$ e $t \in T^*$, de modo que $(S^*, T^*)$ é um corte válido. Para todo $(u, v) \in A$ com $u \in S^*$ e $v \in T^*$, deve-se ter $c_f(u, v) = 0$, implicando $f(u, v) = c(u, v)$. De forma análoga, se $u \in T^*$ e $v \in S^*$, tem-se $f(u, v) = 0$ (caso contrário, existiria capacidade residual reversa $c_f(v, u) > 0$, contradizendo a inalcançabilidade de $u$). O fluxo líquido através do corte satisfaz $|f| = c(S^*, T^*)$, certificando a otimalidade de $f$ e a minimalidade de $(S^*, T^*)$. \hfill $\blacksquare$

\begin{figure}[!ht]
\centering
\begin{tikzpicture}[thick]
    \draw[cut S] (-0.8, -1.8) rectangle (2.8, 1.8);
    \node[cut label S] at (-0.6, 1.65) {$\mathbf{S}$};

    \draw[cut T] (5.4, -1.8) rectangle (9.0, 1.8);
    \node[cut label T] at (8.8, 1.65) {$\mathbf{T}$};

    \node[source node] (s) at (0.0, 0.0) {$s$};
    \node[s node internal] (a) at (2.0, 0.8) {$a$};
    \node[s node internal] (b) at (2.0, -0.8) {$b$};

    \node[t node internal] (c) at (6.2, 0.8) {$c$};
    \node[t node internal] (d) at (6.2, -0.8) {$d$};
    \node[sink node] (t) at (8.2, 0.0) {$t$};

    \draw[line width=1.5pt, dashed, cutviolet] (4.1, -1.85) -- (4.1, 1.85)
        node[above=2.5pt, solid, font=\scriptsize\bfseries, text=cutviolet!90!black, fill=white, inner sep=2.5pt, rounded corners=3pt, draw=cutviolet!70] {Corte Mínimo $(S,T)$};

    \draw[sat edge] (a) -- node[edge lbl above, text=cutviolet!90!black] {$f=c$} (c);
    \draw[sat edge] (b) -- node[edge lbl below, text=cutviolet!90!black] {$f=c$} (d);

    \draw[flow edge] (s) -- (a);
    \draw[flow edge] (s) -- (b);
    \draw[flow edge] (c) -- (t);
    \draw[flow edge] (d) -- (t);
    \draw[flow edge] (a) -- (b);
    \draw[flow edge] (c) -- (d);
\end{tikzpicture}

\caption{Estrutura do corte mínimo $(S,T)$: arcos que cruzam de $S$ para $T$ (em violeta) estão saturados ($f = c$), e fluxos no sentido reverso de $T$ para $S$ são nulos.}
\label{fig:corte_minimo}
\end{figure}

\begin{prop}[Teorema da Integralidade]\label{prop:integralidade_rel}
Se a capacidade $c(a)$ for inteira para todo $a \in A$, então existe um fluxo máximo $f^*$ que é puramente inteiro, isto é, $f^*(a) \in \mathbb{Z}_{\ge 0}$ para todo arco $a \in A$.
\end{prop}

A preservação da integralidade sob capacidades inteiras é garantida pelo fato de que todo aumento ao longo de caminhos residuais em redes inteiras transporta um gargalo $\delta \in \mathbb{Z}_{> 0}$, viabilizando a modelagem exata de problemas combinatórios discretos.

\subsection{Algoritmos de Caminhos Aumentantes}

O algoritmo genérico de Ford-Fulkerson busca caminhos aumentantes via busca em profundidade (DFS). Sob capacidades inteiras, sua complexidade assintótica é $\mathcal{O}(|A| \cdot |f^*|)$, que depende do valor numérico das capacidades (pseudopolinomial). Para obter complexidade polinomial estrita, a literatura introduziu políticas estruturadas de seleção de caminhos.

A Figura~\ref{fig:ford_fulkerson_patologico} ilustra a patologia clássica do método ingênuo com DFS: capacidades de grande magnitude ($M = 10^6$) nos arcos externos e uma conexão transversal unitária ($c=1$) provocam alternâncias sucessivas de aumento unitário ($\delta = 1$), exigindo $2 \times 10^6$ passos. Em contrapartida, algoritmos guiados por distâncias mínimas em arcos (BFS) priorizam caminhos diretos de comprimento 2 ($s \to u \to t$ e $s \to v \to t$), saturando a rede em 2 passos (Edmonds-Karp) ou em uma única fase de fluxo bloqueador (Dinic), desconsiderando o arco transversal por conectar nós no mesmo nível ($d(u) = d(v) = 1$).

\begin{figure}[!ht]
\centering
\resizebox{0.92\textwidth}{!}{%
\begin{tikzpicture}[thick]
\begin{scope}[xshift=0cm]
  \tikzsubcaption{(2.5,2.6)}{(a) Rede patológica ($M = 10^6$)}
  \node[source node] (s) at (0, 1) {$s$};
  \node[flow node] (u) at (2.5, 2) {$u$};
  \node[flow node] (v) at (2.5, 0) {$v$};
  \node[sink node] (t) at (5, 1) {$t$};

  \draw[flow edge] (s) -- node[edge lbl above] {$M$} (u);
  \draw[flow edge] (s) -- node[edge lbl below] {$M$} (v);
  \draw[sat edge] (u) -- node[edge lbl right, text=red!80!black] {$c=1$} (v);
  \draw[flow edge] (u) -- node[edge lbl above] {$M$} (t);
  \draw[flow edge] (v) -- node[edge lbl below] {$M$} (t);
\end{scope}

\begin{scope}[xshift=7.2cm]
  \tikzsubcaption{(2.5,2.6)}{(b) Degeneração da DFS: $2M$ iterações}
  \node[source node] (s2) at (0, 1) {$s$};
  \node[flow node] (u2) at (2.5, 2) {$u$};
  \node[flow node] (v2) at (2.5, 0) {$v$};
  \node[sink node] (t2) at (5, 1) {$t$};

  \draw[tree edge, line width=1.8pt] (s2) -- node[edge lbl above] {$+1$} (u2);
  \draw[sat edge, line width=1.8pt] (u2) -- node[edge lbl right] {$+1$} (v2);
  \draw[tree edge, line width=1.8pt] (v2) -- node[edge lbl below] {$+1$} (t2);

  \draw[flow edge, opacity=0.3] (s2) -- (v2);
  \draw[flow edge, opacity=0.3] (u2) -- (t2);

  \node[font=\scriptsize, align=center, below=3pt] at (2.5,-0.5) {DFS alterna $s \to u \to v \to t$ e $s \to v \to u \to t$\\Aumento $\delta=1 \implies 2 \times 10^6$ passos (BFS / Dinic: 1 fase / 2 passos)};
\end{scope}
\end{tikzpicture}
%
}
\caption{Contraexemplo patológico de Ford-Fulkerson com $M = 10^6$ e arco central unitário $c=1$: (a) rede com fluxo ótimo $|f^*| = 2M$; (b) busca ingênua por DFS alternando caminhos em zigue-zague com aumento $\delta = 1$, demandando $2 \times 10^6$ iterações. Em contraste, métodos baseados em BFS (Edmonds-Karp e Dinic) priorizam caminhos mais curtos ($s \to u \to t$ e $s \to v \to t$), saturando a rede em 2 passos ou em 1 única fase.}
\label{fig:ford_fulkerson_patologico}
\end{figure}

\paragraph{Algoritmo de Edmonds-Karp (1972).} O método seleciona a cada iteração o caminho aumentante de menor quantidade de arcos através de Busca em Largura (BFS) em $D_f$~\cite{edmonds1972theoretical}. Demonstra-se que a distância residual mínima $d(s, v)$ é monotonicamente não decrescente ao longo das iterações. Entre duas saturações consecutivas de um mesmo arco $(u, v)$, a distância da origem sofre um acréscimo de no mínimo duas unidades ($d'(s, u) \ge d(s, u) + 2$). Como a distância máxima em um caminho simples é $|V|-1$, cada arco é saturado no máximo $|V|/2$ vezes. Sendo $|A|$ o número de arcos, ocorrem no máximo $\mathcal{O}(|V| \cdot |A|)$ aumentos. Com custo $\mathcal{O}(|A|)$ por BFS, a complexidade é $\mathcal{O}(|V| \cdot |A|^2)$~\cite{korte2018combinatorial}.

\paragraph{Algoritmo de Dinic (1970).} Yefim Dinitz estruturou o processamento em fases~\cite{dinic1970algorithm}. No início de cada fase, constrói-se o \textbf{Digrafo de Níveis} através de uma BFS global a partir de $s$, retendo apenas os arcos residuais que avançam de nível ($d(v) = d(u) + 1$). Em seguida, emprega-se uma DFS para saturar múltiplos caminhos em um único \textbf{Fluxo Bloqueador} (\textit{Blocking Flow}). Utilizando um vetor de ponteiros mortos (\textit{dead-end pointers}) que evita reexplorar arcos sem capacidade, a fase completa opera em $\mathcal{O}(|V| \cdot |A|)$. Como $d(s, t)$ cresce estritamente a cada fase, ocorrem no máximo $|V|-1$ fases, resultando em complexidade $\mathcal{O}(|V|^2 \cdot |A|)$. Em redes unitárias (como em grafos bipartidos), o número de fases limita-se a $\mathcal{O}(\sqrt{|V|})$, atingindo $\mathcal{O}(|A|\sqrt{|V|})$~\cite{ahuja1993network}.

\begin{figure}[!ht]
\centering
\begin{tikzpicture}[thick]
\foreach \x in {0,2.5,5.0,7.5} { \draw[layer line] (\x,-0.7) -- (\x,3.8); }

\node[gray,font=\footnotesize] at (0,4.1) {nív. 0};
\node[gray,font=\footnotesize] at (2.5,4.1) {nív. 1};
\node[gray,font=\footnotesize] at (5.0,4.1) {nív. 2};
\node[gray,font=\footnotesize] at (7.5,4.1) {nív. 3};

\node[source node] (s) at (0,1.5) {$s$};
\node[flow node] (a) at (2.5,3) {$a$};
\node[flow node] (b) at (2.5,0) {$b$};
\node[flow node] (c) at (5.0,3) {$c$};
\node[flow node] (d) at (5.0,0) {$d$};
\node[sink node] (t) at (7.5,1.5) {$t$};

\draw[flow edge,blue,line width=1.5pt] (s) -- (a);
\draw[flow edge,blue,line width=1.5pt] (s) -- (b);
\draw[flow edge,blue,line width=1.5pt] (a) -- (c);
\draw[flow edge,blue,line width=1.5pt] (b) -- (d);
\draw[flow edge,blue,line width=1.5pt] (a) -- (d);
\draw[flow edge,blue,line width=1.5pt] (b) -- (c);
\draw[flow edge,blue,line width=1.5pt] (c) -- (t);
\draw[flow edge,blue,line width=1.5pt] (d) -- (t);

\draw[nontree edge,bend right=30] (a) to (b);
\draw[nontree edge,bend right=20] (c) to (a);
\end{tikzpicture}

\caption{Digrafo de Níveis no algoritmo de Dinic: arcos azuis são admissíveis ($d(v) = d(u)+1$) e integram a fase de fluxo bloqueador; arcos cinza tracejados não avançam nível e são desconsiderados na fase.}
\label{fig:grafo_niveis}
\end{figure}

\subsection{Algoritmos de Pré-fluxo e Alturas (Push-Relabel)}

A família de algoritmos de Goldberg e Tarjan~\cite{goldberg1988new} adota uma perspectiva dual: relaxa-se a restrição de conservação de fluxo durante o cálculo, mantendo a ausência de caminhos aumentantes como invariante contínua. Um pré-fluxo $f$ respeita as capacidades ($0 \le f(a) \le c(a)$) e admite excessos nodais não negativos $e(v) \ge 0$ para todo $v \in V \setminus \{s, t\}$:
\begin{equation}
    e(v) = \sum_{\{u \mid (u, v) \in A\}} f(u, v) - \sum_{\{w \mid (v, w) \in A\}} f(v, w) \ge 0
\end{equation}

Um vértice com $e(v) > 0$ é denominado \textbf{ativo}. Para direcionar os excessos até o sumidouro $t$ (ou devolvê-los à fonte $s$), mantém-se uma função de alturas $h: V \to \mathbb{N}$ válida, que satisfaz: (i) $h(t) = 0$; (ii) $h(s) = |V|$; e (iii) $h(u) \le h(v) + 1$ para todo $(u, v) \in A_f$. A terceira condição impede a existência de caminhos de $s$ para $t$ em $D_f$. O algoritmo executa duas operações elementares sobre vértices ativos:
\begin{itemize}
    \item \textbf{Push (Empurrar):} Aplicado quando existe arco residual admissível ($c_f(u, v) > 0$ e $h(u) = h(v) + 1$). Envia-se $\delta = \min(e(u), c_f(u, v))$ unidades. Se $\delta = c_f(u, v)$, o push é saturante; caso contrário, é não saturante, zerando $e(u)$;
    \item \textbf{Relabel (Rerrotular):} Aplicado quando $u$ está ativo e não possui vizinhos residuais em nível inferior ($c_f(u, v) > 0 \implies h(u) \le h(v)$). Atualiza-se a altura para $h(u) = 1 + \min \{h(v) \mid c_f(u, v) > 0\}$.
\end{itemize}

Com política de seleção de vértices ativos em fila FIFO, a complexidade assintótica é $\mathcal{O}(|V|^3)$~\cite{goldberg1988new}.

\paragraph{Heurística de Lacuna (\textit{Gap Heuristic}).} Mantendo um vetor de contagem do número de vértices em cada estrato de altura $count[h]$, detecta-se se algum nível $k < |V|$ esvaziou-se ($count[k] = 0$). Como a condição $h(u) \le h(v) + 1$ limita a inclinação dos arcos residuais, nenhum vértice com $h(u) > k$ consegue alcançar o sumidouro ($h(t) = 0$). O algoritmo eleva suas alturas para $|V| + 1$, eliminando operações redundantes de relabel e devolvendo o fluxo diretamente à fonte. Associada à seleção por Maior Altura (\textit{Highest Label}), atinge-se a complexidade $\mathcal{O}(|V|^2 \sqrt{|A|})$~\cite{goldberg1988new}.

\begin{figure}[!ht]
\centering
\begin{tikzpicture}[thick]
  \draw[layer line] (-2, 0) -- (6.5, 0) node[right,black,font=\footnotesize] {$h=0$ (sumidouro $t$)};
  \foreach \y/\h in {1.3/1, 2.6/2, 4.8/4} {
    \draw[layer line] (-2, \y) -- (6.5, \y) node[right,black,font=\footnotesize] {$h=\h$};
  }
  \draw[red!50,dashed,line width=1.5pt] (-2, 3.7) -- (6.5, 3.7) node[right,red!80!black,font=\footnotesize] {$h=3$ (\textbf{LACUNA})};

  \node[sink node,font=\footnotesize] (t) at (2, 0) {$t$};
  \node[flow node,fill=yellow!20,font=\footnotesize] (a) at (2, 1.3) {$a$};
  \node[flow node,fill=yellow!20,font=\footnotesize] (b) at (0.5, 2.6) {$b$};
  \node[flow node,fill=yellow!20,font=\footnotesize] (c) at (3.5, 2.6) {$c$};

  \node[gap node,fill=red!20,font=\footnotesize] (u) at (2, 4.8) {$u$};

  \draw[->,blue!80!black] (b) -- (a);
  \draw[->,blue!80!black] (c) -- (a);
  \draw[->,blue!80!black] (a) -- (t);

  \draw[->,red,line width=1.4pt] (u) -- node[left,pos=0.5,font=\scriptsize,text=black,align=center,xshift=-5mm,fill=white,rounded corners=2pt,inner sep=2pt] {Queda de 2 níveis\\(inadmissível em $A_f$)} (b);
  \draw[->,red,line width=1.4pt] (u) -- node[right,pos=0.5,font=\scriptsize,text=black,align=center,xshift=5mm,fill=white,rounded corners=2pt,inner sep=2pt] {$h(u) \not\le h(c) + 1$\\(inadmissível em $A_f$)} (c);
\end{tikzpicture}

\caption{Mecanismo da Heurística de Lacuna: a restrição $h(u) \le h(v) + 1$ impede arcos residuais com salto de altura superior a uma unidade. O esvaziamento do nível intermediário $h=3$ comprova que vértices com $h \ge 4$ estão desconectados do sumidouro.}
\label{fig:gap_prova}
\end{figure}

\subsection{Reduções Combinatórias em Grafos Bipartidos}

A formulação de fluxo máximo permite resolver problemas clássicos em grafos bipartidos $G = (X \cup Y, E)$ através de reduções polinomiais exatas:

\begin{enumerate}[leftmargin=1.5em, itemsep=0.15em, topsep=0.2em, parsep=0em]
    \item \textbf{Emparelhamento Máximo Bipartido (\textit{Maximum Bipartite Matching}):} Constrói-se uma rede direcionada adicionando uma superfonte $s$ conectada a cada $x \in X$ com capacidade 1, e cada $y \in Y$ conectado ao supersumidouro $t$ com capacidade 1. Cada aresta $\{x, y\} \in E$ torna-se um arco $(x, y)$ com capacidade infinita (ou unitária). Pelo Teorema da Integralidade, o fluxo máximo tem valor inteiro e coincide rigorosamente com a cardinalidade máxima de emparelhamento: $|f^*| = \nu(G)$.

\begin{figure}[!ht]
\centering
\begin{tikzpicture}[thick]
\begin{scope}[xshift=0cm]
  \tikzsubcaption{(1.5,3.8)}{(a) Grafo bipartido original}
  \foreach \i/\y in {1/3, 2/1.5, 3/0} {
    \node[bip L] (l\i) at (0,\y) {$l_\i$};
    \node[bip R] (r\i) at (3,\y) {$r_\i$};
  }
  \draw[-] (l1) -- (r1); \draw[-] (l1) -- (r2);
  \draw[-] (l2) -- (r2); \draw[-] (l2) -- (r3);
  \draw[-] (l3) -- (r1); \draw[-] (l3) -- (r3);
\end{scope}
\node at (4.8,1.5) {\Huge $\Rightarrow$};

\begin{scope}[xshift=6.2cm]
  \tikzsubcaption{(2.5,3.8)}{(b) Rede de fluxo equivalente $D'$}
  \node[source node] (s) at (0,1.5) {$s$};
  \node[sink node] (t) at (5,1.5) {$t$};
  \foreach \i/\y in {1/3, 2/1.5, 3/0} {
    \node[bip L,font=\footnotesize] (fl\i) at (1.8,\y) {$l_\i$};
    \node[bip R,font=\footnotesize] (fr\i) at (3.2,\y) {$r_\i$};
    \draw[flow edge] (s) -- node[above,font=\footnotesize] {$1$} (fl\i);
    \draw[flow edge] (fr\i) -- node[above,font=\footnotesize] {$1$} (t);
  }
  \draw[flow edge,blue,line width=1.5pt] (fl1) -- (fr1); \draw[flow edge] (fl1) -- (fr2);
  \draw[flow edge,blue,line width=1.5pt] (fl2) -- (fr2); \draw[flow edge] (fl2) -- (fr3);
  \draw[flow edge] (fl3) -- (fr1); \draw[flow edge,blue,line width=1.5pt] (fl3) -- (fr3);
\end{scope}
\end{tikzpicture}

\caption{Redução do problema de Emparelhamento Máximo Bipartido para Fluxo Máximo: (a) grafo bipartido original; (b) rede com superfonte $s$, supersumidouro $t$ e capacidades unitárias. Os arcos azuis com $f(u, v) = 1$ indicam as arestas do emparelhamento máximo $|f^*| = 3$.}
\label{fig:fluxo_bipartido}
\end{figure}

    \item \textbf{Teorema de König (1931) e Cobertura Mínima de Vértices:} Em todo grafo bipartido, a cardinalidade do emparelhamento máximo é igual à cardinalidade da cobertura mínima de vértices ($\tau(G)$):
    \begin{equation}
        \nu(G) = \tau(G)
    \end{equation}
    A demonstração constrói o corte mínimo residual $(S^*, T^*)$ em $D_{f^*}$ a partir de $s$. O conjunto $C^* = (X \setminus S^*) \cup (Y \cap S^*)$ cobre todas as arestas do grafo e sua cardinalidade satisfaz $|C^*| = c(S^*, T^*) = |f^*| = \nu(G)$, fornecendo um algoritmo polinomial em tempo $\mathcal{O}(|E|\sqrt{|V|})$ via Dinic.

    \item \textbf{Conjunto Independente Máximo e Identidades de Gallai (1959):} Um subconjunto de vértices $I$ é independente se e somente se seu complementar $V \setminus I$ for uma cobertura de vértices. Pelo Teorema de Gallai, $\alpha(G) + \tau(G) = |V|$ e $\nu(G) + \rho(G) = |V|$, onde $\alpha(G)$ é a cardinalidade do conjunto independente máximo e $\rho(G)$ da cobertura mínima de arestas. Em grafos bipartidos, $\alpha(G) = \rho(G) = |V| - \nu(G)$, e o conjunto ótimo é extraído diretamente por $I^* = V \setminus C^*$ em tempo $\mathcal{O}(|V|)$.

    \item \textbf{Teorema do Casamento de Hall (1935):} Um grafo bipartido $G = (X \cup Y, E)$ possui um emparelhamento que cobre todos os vértices de $X$ se e somente se $|N(S)| \ge |S|$ para todo $S \subseteq X$, onde $N(S)$ é a vizinhança de $S$. A prova deduz-se diretamente de König aplicando a contradição de corte mínimo.
\end{enumerate}

\subsection{Conectividade e Roteamento}

A determinação do número de rotas independentes entre dois vértices $s$ e $t$ reduz-se a fluxo máximo:
\begin{itemize}[leftmargin=1.5em, itemsep=0.15em, topsep=0.2em, parsep=0em]
    \item \textbf{Caminhos Disjuntos por Arcos:} Fixando $c(a) = 1$ para todo arco, o valor $|f^*|$ equivale ao número máximo de rotas direcionadas sem arcos em comum (Teorema de Menger por arcos, 1927).
    \item \textbf{Caminhos Disjuntos por Vértices e \textit{Vertex Splitting}:} Para impedir que rotas compartilhem vértices intermediários, divide-se cada vértice $v \in V \setminus \{s, t\}$ em dois nós, $v_{\text{in}}$ e $v_{\text{out}}$, conectados por um arco interno $(v_{\text{in}}, v_{\text{out}})$ com capacidade unitária $c(v_{\text{in}}, v_{\text{out}}) = 1$. O corte mínimo nesta rede expandida identifica o conjunto de vértices cuja remoção desconecta a fonte do sumidouro (Teorema de Menger por vértices, 1927).
\end{itemize}

\subsection{Teoria do Fluxo de Custo Mínimo (MCF)}

No problema de Fluxo de Custo Mínimo, associa-se a cada arco $(u, v) \in A$ uma capacidade $c(u, v) \ge 0$ e um custo linear unitário $w(u, v) \in \mathbb{R}$. Cada vértice $v \in V$ possui demanda $b(v) \in \mathbb{R}$ com $\sum_{v \in V} b(v) = 0$. O problema formula-se pelo programa linear:
\begin{align}
    \min_{f} \quad & \sum_{(u, v) \in A} w(u, v) \cdot f(u, v) \\
    \sa \quad & \sum_{\{w \mid (v, w) \in A\}} f(v, w) - \sum_{\{u \mid (u, v) \in A\}} f(u, v) = b(v), \quad \forall v \in V \\
    & 0 \le f(u, v) \le c(u, v), \quad \forall (u, v) \in A
\end{align}

Na rede residual $D_f$, os arcos diretos preservam o custo $w_f(u, v) = w(u, v)$ e os arcos reversos recebem custo antissimétrico $w_f(v, u) = -w(u, v)$, refletindo a devolução de custo ao cancelar fluxo.

\begin{prop}[Condição de Otimalidade por Ciclos Negativos --- Klein, 1967]\label{prop:klein_rel}
Um fluxo viável $f$ é de custo mínimo se, e somente se, o digrafo residual associado $D_f$ não contém nenhum ciclo direcionado de custo total estritamente negativo.
\end{prop}

\begin{prop}[Custos Reduzidos e Potenciais Nodais]\label{prop:potenciais_rel}
Seja $\pi: V \to \mathbb{R}$ um vetor de potenciais nodais (variáveis duais das restrições de balanço). O custo reduzido de um arco residual $(u, v) \in A_f$ é definido por:
\begin{equation}
    \bar{w}_\pi(u, v) = w_f(u, v) - \pi(u) + \pi(v)
\end{equation}
Um fluxo viável $f$ é ótimo se e somente se existe $\pi$ tal que $\bar{w}_\pi(u, v) \ge 0$ para todo $(u, v) \in A_f$.
\end{prop}

A não negatividade dos custos reduzidos corresponde às condições de Folga Complementar de Karush-Kuhn-Tucker (KKT). Mantendo os potenciais atualizados ao longo das iterações, viabiliza-se o uso do algoritmo de Dijkstra com fila de prioridades no lugar de algoritmos com suporte a custos negativos no método \textit{Successive Shortest Path} (SSP).

\subsection{O Algoritmo Network Simplex}

O algoritmo \textit{Network Simplex} especializa o método Simplex linear explorando a estrutura combinatória de árvores geradoras da rede (veja~\cite{ahuja1993network, dantzig1951application}). Uma solução básica viável particiona os arcos em três classes $(T, L, U)$:
\begin{enumerate}
    \item $T$: arcos básicos que formam uma árvore geradora do digrafo;
    \item $L$: arcos não básicos com fluxo no limite inferior ($f(a) = 0$);
    \item $U$: arcos não básicos com fluxo saturado no limite superior ($f(a) = c(a)$).
\end{enumerate}

\begin{prop}[Total Unimodularidade (TUM) da Matriz Nó-Arco]\label{prop:tum_rel}
Seja $\mathbf{M} \in \{-1, 0, 1\}^{|V| \times |A|}$ a matriz de incidência nó-arco de $D$. Todo determinante de submatriz quadrada de $\mathbf{M}$ pertence a $\{-1, 0, 1\}$. Consequentemente, para qualquer vetor de demandas e capacidades inteiras, toda solução básica (base em árvore) é puramente inteira.
\end{prop}

A Total Unimodularidade garante que a matriz básica $\mathbf{B}$ associada a qualquer árvore $T$ satisfaça $|\det \mathbf{B}| = 1$, eliminando divisões na Regra de Cramer e garantindo aritmética inteira exata.

Cada iteração do Network Simplex executa quatro passos: (i) \textbf{Pricing:} identificação de arco violador $(u, v)$ fora da árvore ($\bar{w}_\pi < 0$ em $L$ ou $\bar{w}_\pi > 0$ em $U$); (ii) \textbf{Ciclo Fundamental:} determinação do ciclo $C$ formado por $(u, v)$ em $T$ via busca do ancestral comum mais próximo (\textit{LCA}); (iii) \textbf{Gargalo:} cômputo da capacidade residual mínima $\delta$ em $C$ e determinação do arco que deixa a base; e (iv) \textbf{Pivoteamento:} atualização do fluxo e ajuste dos potenciais nodais na subárvore em tempo $\mathcal{O}(|V|)$ com auxílio de listas de pré-ordem (\lstinline{thread}).

\begin{figure}[!ht]
\centering
\begin{tikzpicture}[thick]
\begin{scope}[xshift=0cm]
  \tikzsubcaption{(2.0,3.6)}{(a) Entrada do arco $(3,4)$ no ciclo fundamental}
  \node[source node] (a1) at (0, 1.3) {$1$};
  \node[flow node] (a2) at (2, 2.6) {$2$};
  \node[flow node] (a3) at (2, 0.0) {$3$};
  \node[sink node] (a4) at (4, 1.3) {$4$};

  \draw[tree edge] (a1) -- (a2);
  \draw[tree edge] (a2) -- (a4);
  \draw[tree edge] (a1) -- (a3);
  \draw[pivot enter] (a3) -- node[below=3pt,sloped,font=\scriptsize,text=red!80!black] {entra} (a4);
\end{scope}

\node at (5.2,1.3) {\Large $\Rightarrow$};

\begin{scope}[xshift=6.4cm]
  \tikzsubcaption{(2.0,3.6)}{(b) Nova base em árvore após o pivot}
  \node[source node] (b1) at (0, 1.3) {$1$};
  \node[flow node] (b2) at (2, 2.6) {$2$};
  \node[flow node] (b3) at (2, 0.0) {$3$};
  \node[sink node] (b4) at (4, 1.3) {$4$};

  \draw[tree edge] (b1) -- (b2);
  \draw[pivot leave] (b2) -- node[above=3pt,sloped,font=\scriptsize,text=gray!80!black] {sai} (b4);
  \draw[tree edge] (b1) -- (b3);
  \draw[tree edge] (b3) -- (b4);
\end{scope}
\end{tikzpicture}

\caption{Iteração do Network Simplex: (a) o arco não básico violador $(3,4)$ com $\bar{w}_\pi < 0$ fecha um ciclo fundamental em $T$; (b) o fluxo é aumentado por $\delta$ e o arco saturado $(2,4)$ deixa a base, restabelecendo a base em árvore.}
\label{fig:simplex_pivoteamento}
\end{figure}

% =========================================================================
% SEÇÃO 4: RESULTADOS COMPUTACIONAIS
% =========================================================================
\section{Resultados Computacionais}\label{sec:resultados}

Esta seção apresenta a avaliação experimental sistemática dos nove motores algorítmicos desenvolvidos em C++23, submetidos a testes de estresse em instâncias canônicas da literatura acadêmica.

\subsection{Protocolo de Teste e Isolamento de CPU}

Para garantir rigor científico e reprodutibilidade, adotaram-se os seguintes controles metodológicos:
\begin{enumerate}[leftmargin=1.5em, itemsep=0.15em, topsep=0.2em, parsep=0em]
    \item \textbf{Isolamento Estrito de CPU:} A medição temporal utiliza o relógio monotônico de hardware \lstinline{std::chrono::high_resolution_clock}. As etapas de I/O em disco, leitura e parsing de arquivos DIMACS e alocação inicial de memória são estritamente excluídas do cronômetro; mede-se exclusivamente o método de cálculo;
    \item \textbf{Regime Estacionário (\textit{Steady State}):} Cada instância é executada de forma independente por $N = 5$ repetições sucessivas após aquecimento de cache, reportando-se a média amostral ($\bar{t}$ em milissegundos) e o desvio-padrão;
    \item \textbf{Validação Cruzada de Otimalidade:} Como os algoritmos operam sob paradigmas distintos (aumentos primais, pré-fluxos duais e pivoteamento em árvore), a suíte verifica automaticamente se todos os motores convergem exatamente para a mesma vazão máxima $f^*$ e custo mínimo $z^*$. Em 100\% dos testes, houve concordância exata de soluções ótimas.
\end{enumerate}

\subsection{Avaliação de Fluxo Máximo (Coleções DIMACS)}

Os motores de fluxo máximo foram avaliados sobre as coleções clássicas do \textit{First DIMACS Challenge}:
\begin{itemize}[leftmargin=1.5em, itemsep=0.15em, topsep=0.2em, parsep=0em]
    \item \textbf{Washington (RLG e Line):} Redes com gargalos induzidos para provocar caminhos aumentantes longos;
    \item \textbf{Genrmf (small, medium, wide, long):} Grafos em grade tridimensional com planos transversais de corte severo, representando o pior cenário para métodos de pré-fluxo e caminhos aumentantes.
\end{itemize}

A Tabela~\ref{tab:benchmarks_maxflow} apresenta os tempos de execução médios obtidos pelos cinco motores nas instâncias de referência de fluxo máximo.

\begin{table}[!ht]
\centering
\small
\setlength{\tabcolsep}{3.5pt}
\begin{tabular}{lrrrrrrrr}
\toprule
\textbf{Instância} & $|V|$ & $|A|$ & $f^*$ & \textbf{FF (ms)} & \textbf{EK (ms)} & \textbf{Dinic (ms)} & \textbf{PR (ms)} & \textbf{PR-Gap (ms)} \\
\midrule
\texttt{wash\_rlg\_16}  & 18    & 48    & 87    & 0,004 & 0,006 & \textbf{0,003} & 0,008 & \textbf{0,003} \\
\texttt{wash\_rlg\_32}  & 34    & 96    & 134   & 0,012 & 0,015 & \textbf{0,006} & 0,019 & 0,007 \\
\texttt{wash\_rlg\_64}  & 66    & 192   & 198   & 0,035 & 0,042 & \textbf{0,014} & 0,051 & 0,018 \\
\texttt{wash\_line\_16} & 18    & 32    & 1     & 0,003 & 0,004 & \textbf{0,002} & 0,005 & \textbf{0,002} \\
\texttt{wash\_line\_64} & 66    & 128   & 1     & 0,015 & 0,022 & \textbf{0,007} & 0,031 & 0,009 \\
\texttt{genrmf\_small}  & 256   & 1.056 & 1.450 & 1,420 & 1,890 & 0,310 & 0,840 & \textbf{0,180} \\
\texttt{genrmf\_med}    & 1.024 & 4.608 & 2.890 & 12,80 & 18,40 & 1,950 & 6,200 & \textbf{0,920} \\
\texttt{genrmf\_wide}   & 4.096 & 19.456& 5.120 & 98,50 & 142,1 & 14,20 & 42,80 & \textbf{5,600} \\
\texttt{genrmf\_long}   & 4.096 & 18.432& 1.820 & 185,2 & 280,4 & 18,90 & 68,50 & \textbf{7,100} \\
\bottomrule
\end{tabular}
\caption{Comparação experimental dos motores de fluxo máximo em instâncias DIMACS.}
\label{tab:benchmarks_maxflow}
\end{table}

\noindent \textbf{Análise dos Resultados de Fluxo Máximo:}
\begin{itemize}[leftmargin=1.5em, itemsep=0.15em, topsep=0.2em, parsep=0em]
    \item \textbf{Impacto da Gap Heuristic no Push-Relabel:} O motor \code{PushRelabelImproved} superou o \code{PushRelabel} convencional em todas as instâncias médias e grandes, atingindo uma aceleração de até \textbf{$9{,}7\times$} na instância \texttt{genrmf\_long} (7,10 ms vs. 68,50 ms), comprovando a eficácia da detecção precoce de rótulos mortos;
    \item \textbf{Dinic vs. Edmonds-Karp:} O algoritmo de Dinic foi de \textbf{$6\times$ a $15\times$ mais rápido} que o Edmonds-Karp em redes de grande escala. Enquanto o Edmonds-Karp executa uma nova BFS global a cada aumento de fluxo ($O(n m^2)$), o Dinic constrói o digrafo de níveis uma única vez por fase e empurra múltiplos caminhos em fluxo bloqueador com eliminação de ponteiros mortos;
    \item \textbf{Degradação sob Gargalos Profundos:} Em redes lineares e RMF longas, Ford-Fulkerson e Edmonds-Karp degradam acentuadamente em função do número de aumentos e do comprimento médio dos caminhos, enquanto os algoritmos baseados em alturas mantêm excelente estabilidade.
\end{itemize}

\subsection{Avaliação de Fluxo de Custo Mínimo (Coleção DIMACS Netgen)}

Para avaliar os motores de fluxo de custo mínimo, utilizou-se a família \textbf{Netgen}, gerada oficialmente pelo gerador DIMACS de Rutgers para redes de transporte com restrições balanceadas de oferta e demanda.

A Tabela~\ref{tab:benchmarks_mincost} apresenta a comparação entre os quatro motores de custo mínimo nas instâncias \texttt{netgen\_01} a \texttt{netgen\_08}.

\begin{table}[!ht]
\centering
\small
\setlength{\tabcolsep}{3.5pt}
\begin{tabular}{lrrrrrrrr}
\toprule
\textbf{Instância} & $|V|$ & $|A|$ & Fluxo ($f^*$) & Custo ($z^*$) & \textbf{CC (ms)} & \textbf{SSP-SPFA} & \textbf{SSP-Dijk} & \textbf{Simplex (ms)} \\
\midrule
\texttt{netgen\_01} & 100 & 500   & 5.000  & 142.800 & 42,10 & 18,50 & 5,20 & \textbf{0,62} \\
\texttt{netgen\_02} & 150 & 800   & 7.500  & 289.400 & 128,4 & 45,20 & 11,80 & \textbf{1,15} \\
\texttt{netgen\_03} & 200 & 1.200 & 10.000 & 481.200 & 310,5 & 92,10 & 22,40 & \textbf{1,80} \\
\texttt{netgen\_04} & 250 & 1.500 & 12.500 & 672.100 & 580,2 & 165,4 & 38,10 & \textbf{2,45} \\
\texttt{netgen\_05} & 300 & 2.000 & 15.000 & 915.000 & 1.120 & 280,3 & 59,20 & \textbf{3,50} \\
\texttt{netgen\_06} & 400 & 3.000 & 20.000 & 1.450.200 & 2.890 & 540,1 & 112,0 & \textbf{5,80} \\
\texttt{netgen\_07} & 500 & 4.000 & 25.000 & 2.120.400 & 5.420 & 980,5 & 195,4 & \textbf{8,90} \\
\texttt{netgen\_08} & 600 & 5.000 & 30.000 & 2.980.100 & 9.850 & 1.620 & 310,2 & \textbf{12,40} \\
\bottomrule
\end{tabular}
\caption{Comparação experimental dos motores de fluxo de custo mínimo em instâncias Netgen.}
\label{tab:benchmarks_mincost}
\end{table}

\noindent \textbf{Análise dos Resultados de Custo Mínimo:}
\begin{itemize}[leftmargin=1.5em, itemsep=0.15em, topsep=0.2em, parsep=0em]
    \item \textbf{Desempenho do Network Simplex:} O algoritmo \code{NetworkSimplex} apresentou menor tempo de execução que os demais métodos, sendo até \textbf{$25\times$ mais rápido que o SSP com Dijkstra} e aproximadamente \textbf{$800\times$ mais rápido que o Cycle Canceling} na instância \texttt{netgen\_08} (12,40 ms vs. 9.850 ms). A eficiência decorre da estrutura de árvore geradora: enquanto métodos de caminhos mínimos recalculam trajetos completos a cada iteração, o Network Simplex executa pivoteamentos locais em ciclos fundamentais, atualizando potenciais apenas na subárvore afetada;
    \item \textbf{Potenciais Nodais no SSP (Dijkstra vs. SPFA):} O uso de potenciais nodais ($\bar{w}_\pi \ge 0$) reduziu o tempo de execução do \textit{Successive Shortest Path} em \textbf{$3{,}5\times$ a $5{,}2\times$} frente ao SSP com SPFA, confirmando a vantagem de Dijkstra com fila de prioridades;
    \item \textbf{Limitações do Cycle Canceling:} O cancelamento de ciclos negativos requer buscas completas por ciclos via Bellman-Ford ou SPFA ($O(n m)$ por iteração). Em instâncias com muitas unidades de fluxo ou redes com mais de 500 vértices, o método apresenta tempos significativamente superiores aos demais.
\end{itemize}

% =========================================================================
% SEÇÃO 5: APLICAÇÃO EM SEGMENTAÇÃO DE IMAGENS
% =========================================================================
\section{Aplicação Prática: Segmentação Interativa de Imagens}\label{sec:segmentacao}

Esta seção apresenta a modelagem e a avaliação experimental da \textbf{Segmentação Interativa de Imagens} via corte mínimo em grafos (\textit{Graph Cuts}), implementada no módulo \texttt{Aplicações/Segmentação de Imagens/}.

\subsection{Modelagem Energética de Boykov-Jolly}

Seja $I$ uma imagem digital retangular composta por $n = W \times H$ pixels. O objetivo consiste em associar a cada pixel $p \in \mathcal{P}$ um rótulo binário $f_p \in \{0, 1\}$, onde $0$ representa o fundo (\textit{background}) e $1$ o objeto de primeiro plano (\textit{foreground}). A função de energia de Boykov e Jolly~\cite{boykov2001interactive} pondera a fidelidade às características de cor e a continuidade espacial de contornos:
\begin{equation}
    E(f) = \sum_{p \in \mathcal{P}} D_p(f_p) + \sum_{\{p, q\} \in \mathcal{N}} V_{p, q}(f_p, f_q)
\end{equation}
onde:
\begin{itemize}[leftmargin=1.5em, itemsep=0.15em, topsep=0.2em, parsep=0em]
    \item $D_p(f_p)$ é o termo de dados (\textit{Data Term}), que penaliza atribuições de rótulos incoerentes com os perfis de cor das sementes marcadas pelo usuário;
    \item $V_{p, q}(f_p, f_q)$ é o termo de suavidade (\textit{Smoothness Term}) sobre a vizinhança 4-conectada $\mathcal{N}$, atribuindo custo $V(0, 1) = V(1, 0) = w_n(p, q)$ a pixels vizinhos com rótulos opostos e custo 0 se compartilham o mesmo rótulo.
\end{itemize}

\subsection{Mapeamento em Digrafo e Equivalência Exata}

A rede $D = (V, A)$ é construída com $|V| = n + 2$ nós, adicionando uma superfonte $s$ (\textit{foreground}) e um supersumidouro $t$ (\textit{background}). As capacidades dos arcos são estabelecidas por:
\begin{enumerate}[leftmargin=1.5em, itemsep=0.15em, topsep=0.2em, parsep=0em]
    \item \textbf{N-links (Vizinhança 4-Conectada):} Entre cada par de pixels adjacentes $(p, q)$, insere-se uma conexão com capacidade inversamente proporcional ao gradiente cromático no espaço tridimensional RGB:
    \begin{equation}
        w_n(p, q) = \left\lfloor K \cdot \exp\left(-\frac{\|I_p - I_q\|^2}{2\sigma^2}\right) \right\rfloor
    \end{equation}
    onde $\|I_p - I_q\|$ é a distância euclidiana de cor e $\sigma$ regula a sensibilidade a contrastes. Pixels homogêneos recebem peso elevado, direcionando o corte para bordas de variação brusca;
    \item \textbf{T-links (Conexões Terminais):}
    \begin{itemize}
        \item Para sementes de primeiro plano ($\mathcal{S}_{\text{fg}}$): $w_s(p) = \infty$ e $w_t(p) = 0$;
        \item Para sementes de fundo ($\mathcal{S}_{\text{bg}}$): $w_s(p) = 0$ e $w_t(p) = \infty$;
        \item Para pixels não marcados, utiliza-se a distância euclidiana mínima de cor frente às sementes de cada classe:
        \begin{equation}
            w_s(p) = \max\left( \left\lfloor K \cdot \exp\left(-\frac{\min_{s_j \in \mathcal{S}_{\text{fg}}} \|I_p - I_{s_j}\|^2}{2\sigma^2}\right) \right\rfloor, 1 \right)
        \end{equation}
        \begin{equation}
            w_t(p) = \max\left( \left\lfloor K \cdot \exp\left(-\frac{\min_{s_j \in \mathcal{S}_{\text{bg}}} \|I_p - I_{s_j}\|^2}{2\sigma^2}\right) \right\rfloor, 1 \right)
        \end{equation}
    \end{itemize}
\end{enumerate}

\begin{prop}[Equivalência Exata Min-Cut $\equiv$ Min-Energy]
Seja $(S, T)$ um $st$-corte finito na rede de segmentação, induzindo a rotulação $f_p = 1$ se $p \in S$ e $f_p = 0$ se $p \in T$. A capacidade do corte satisfaz:
\begin{equation}
    c(S, T) = E(f) + \text{constante aditiva}
\end{equation}
onde a constante depende unicamente da rede e independe de $f$. Portanto, o corte mínimo obtido por algoritmos de fluxo determina o mínimo global exato de $E(f)$ em tempo polinomial.
\end{prop}

A Figura~\ref{fig:grid_cut_3d} ilustra a estrutura da rede tridimensional gerada pelo modelo.

\begin{figure}[!ht]
\centering
\resizebox{0.75\textwidth}{!}{%
\begin{tikzpicture}[
    x={(-1.2cm, -0.6cm)},
    y={(1.2cm, -0.6cm)},
    z={(0cm, 1.5cm)}
]
    \node[sink node] (t) at (2,2,-3.2) {$t$ (Fundo)};

    \foreach \x in {1,2,3} {
        \foreach \y in {1,2,3} {
            \coordinate (cp\x\y) at (\x,\y,0);
        }
    }

    \foreach \x in {1,2,3} {
        \foreach \y in {1,2,3} {
            \draw[->, red!70!black, dashed, line width=0.8pt, draw opacity=0.65] (cp\x\y) -- (t);
        }
    }

    \draw[fill=gray!15, draw=gray!60, fill opacity=0.85] (0.3,0.3,0) -- (3.7,0.3,0) -- (3.7,3.7,0) -- (0.3,3.7,0) -- cycle;

    \foreach \x in {1,2,3} {
        \foreach \y in {1,2} {
            \pgfmathtruncatemacro{\ynext}{\y+1}
            \draw[<->, line width=1.1pt, black!75] (cp\x\y) -- (cp\x\ynext);
        }
    }
    \foreach \x in {1,2} {
        \foreach \y in {1,2,3} {
            \pgfmathtruncatemacro{\xnext}{\x+1}
            \draw[<->, line width=1.1pt, black!75] (cp\x\y) -- (cp\xnext\y);
        }
    }

    \foreach \x in {1,2,3} {
        \foreach \y in {1,2,3} {
            \node[circle, fill=blue!60!black, inner sep=2.2pt] (p\x\y) at (cp\x\y) {};
        }
    }

    \node[source node] (s) at (2,2,3.2) {$s$ (Objeto)};

    \foreach \x in {1,2,3} {
        \foreach \y in {1,2,3} {
            \draw[->, blue!70!black, line width=0.8pt, draw opacity=0.7] (s) -- (p\x\y);
        }
    }
\end{tikzpicture}%
}
\caption{Estrutura da rede de segmentação: superfonte $s$ e supersumidouro $t$ conectados aos pixels por T-links verticais, com N-links bidirecionais entre pixels adjacentes.}
\label{fig:grid_cut_3d}
\end{figure}

\subsection{Resultados Experimentais da Segmentação}

O pipeline de segmentação foi validado sobre as imagens de teste do repositório (\textit{balloon}, \textit{circle}, \textit{duck}, \textit{grid1}, \textit{grid2}), confrontando as saídas geradas em formato binário PPM com as máscaras gabarito (\textit{ground-truth}) via comparação exata byte a byte (\lstinline{cmp}). Em todas as instâncias, obteve-se 100\% de concordância com o gabarito.

A Tabela~\ref{tab:segmentacao_resultados} apresenta a escalabilidade do pipeline em função das dimensões das imagens, número de vértices, arcos residuais e tempo de processamento do corte mínimo com o algoritmo \code{Dinic}.

\begin{table}[!ht]
\centering
\small
\setlength{\tabcolsep}{5pt}
\begin{tabular}{lrrrrr}
\toprule
\textbf{Instância} & \textbf{Dimensões ($W \times H$)} & \textbf{Vértices ($|V|$)} & \textbf{Arcos ($|A|$)} & \textbf{Custo do Corte ($E^*$)} & \textbf{Tempo de CPU (ms)} \\
\midrule
\texttt{grid1.ppm}    & $10 \times 10$   & 102     & 460       & 1.840       & 0,12 \\
\texttt{circle.ppm}   & $50 \times 50$   & 2.502   & 14.700    & 28.450      & 2,41 \\
\texttt{balloon.ppm}  & $100 \times 100$ & 10.002  & 69.400    & 142.110     & 14,85 \\
\texttt{duck.ppm}     & $150 \times 150$ & 22.502  & 156.900   & 389.200     & 48,20 \\
\texttt{balloons.ppm} & $200 \times 200$ & 40.002  & 279.400   & 715.630     & 95,70 \\
\bottomrule
\end{tabular}
\caption{Desempenho do pipeline de segmentação via corte mínimo sobre instâncias reais.}
\label{tab:segmentacao_resultados}
\end{table}

Mesmo em redes com mais de 40.000 vértices e 270.000 arcos, o corte mínimo global é obtido em menos de 100 milissegundos.

% =========================================================================
% SEÇÃO 6: CONCLUSÕES
% =========================================================================
\section{Conclusões}\label{sec:conclusoes}

Os resultados obtidos demonstram o cumprimento integral dos objetivos do plano de trabalho, articulando teoria matemática, engenharia de software em C++23 e validação experimental:

\begin{enumerate}[leftmargin=1.5em, itemsep=0.2em, topsep=0.2em, parsep=0em]
    \item \textbf{Fundamentação Teórica:} Demonstração formal dos teoremas centrais da teoria de fluxos: Fluxo Máximo e Corte Mínimo, Integralidade de Ford-Fulkerson, Decomposição de Fluxo, Teorema de Menger (arcos e vértices), Teorema de König e as Quatro Dualidades em grafos bipartidos. No domínio de custo mínimo: condição de otimalidade por ciclos negativos de Klein, equivalência dos custos reduzidos sob potenciais nodais, complementaridade de folgas KKT e Total Unimodularidade da matriz de incidência;

    \item \textbf{Arquitetura em C++23:} Desenvolvimento de classes organizadas em duas famílias polimórficas com representação de adjacência contígua e reversão em $\mathcal{O}(1)$:
    \begin{itemize}[leftmargin=1.5em, itemsep=0.1em, topsep=0.1em, parsep=0em]
        \item \textbf{Família \lstinline{FlowNetwork} (Fluxo Máximo):} Centralizada na classe base abstrata \lstinline{FlowNetwork} e implementada pelos algoritmos \lstinline{FordFulkerson}, \lstinline{EdmondsKarp}, \lstinline{Dinic}, \lstinline{PushRelabel} (FIFO) e \lstinline{PushRelabelImproved} (Gap Heuristic), documentados nos Apêndices~\ref{ap:flownetwork} a~\ref{ap:push_relabel_improved};
        \item \textbf{Família \lstinline{CostNetwork} (Custo Mínimo):} Centralizada na classe base abstrata \lstinline{CostNetwork} e implementada pelos algoritmos \lstinline{CycleCanceling}, \lstinline{SuccessiveShortest} (SPFA), \lstinline{SuccessiveShortestDijkstra} (Potenciais) e \lstinline{NetworkSimplex}, documentados nos Apêndices~\ref{ap:costnetwork} a~\ref{ap:network_simplex};
    \end{itemize}

    \item \textbf{Reduções Combinatórias e Aplicações:} Resolução de problemas clássicos (\textit{Maximum Bipartite Matching}, \textit{Minimum Vertex Cover}, \textit{Maximum Independent Set}, caminhos disjuntos e desconexão de redes) e implementação de \textbf{Segmentação Interativa de Imagens} baseada em corte mínimo via corte de grafos (\textit{Graph Cuts});

    \item \textbf{Avaliação Experimental:} Protocolo experimental com instâncias da literatura (DIMACS Washington, Genrmf e Netgen), utilizando isolamento estrito de tempo de CPU e validação cruzada com 100\% de concordância entre os algoritmos implementados.
\end{enumerate}

Como trabalhos futuros, destacam-se a avaliação empírica sobre instâncias esparsas massivas, a análise de diferentes regras de pivoteamento no \textit{Network Simplex} (como \textit{Candidate List Pricing}) e a preparação de artigo científico com a síntese dos resultados.


\newpage
\section{Bibliografia}
\bibliographystyle{plain}
\bibliography{cit}

\newpage
\appendix
\renewcommand{\theHsection}{apendice.\arabic{section}}
\providecommand{\apendice}[1]{\section{#1}}

\begin{center}{\Large\textbf{Apêndice Computacional: Implementações em C++23}}\end{center}

% =========================================================================
% ARQUITETURA DE SOFTWARE EM C++23 E ESTRUTURAS DE DADOS
% =========================================================================
\apendice{Arquitetura Orientada a Objetos em C++23 e Estruturas de Dados}
\label{ap:arquitetura}

A transposição dos modelos matemáticos abstratos para motores computacionais de alto desempenho fundamenta-se no paradigma da Orientação a Objetos através do padrão ISO C++23.

\paragraph{Tipagem e Margem de Segurança Numérica.}
A biblioteca adota apelidos de tipo (\textit{Type Aliases}) canônicos: \lstinline{Long} (\lstinline{long long}) para capacidades e fluxos, e \lstinline{Size} (\lstinline{std::size_t}) para identificadores nodais e dimensionamento de estruturas de dados. Para mitigar o risco de estouro de inteiros (\textit{Integer Overflow}) em somas cumulativas de capacidades infinitas, as constantes de saturação são definidas com deslocamento defensivo de 8 bits à direita:
\begin{lstlisting}[language=C++, numbers=none, basicstyle=\ttfamily\footnotesize, frame=none, backgroundcolor=\color{white}]
constexpr Long INF = std::numeric_limits<Long>::max() >> 8;
constexpr Size MAX = std::numeric_limits<Size>::max() >> 8;
\end{lstlisting}

\paragraph{Representação Contígua de Grafo Residual e Emparelhamento XOR.}
Para favorecer a localidade de cache da CPU e eliminar o overhead de alocações dinâmicas dispersas, a topologia da rede é mantida em duas estruturas integradas:
\begin{enumerate}
    \item \lstinline{std::vector<Edge> edges}: vetor contíguo na memória contendo todas as conexões da rede (\lstinline{from}, \lstinline{to}, \lstinline{capacity}, \lstinline{flow});
    \item \lstinline{std::vector<std::vector<Size>> adj}: lista de adjacência onde cada entrada armazena apenas os índices numéricos dos arcos no vetor principal \lstinline{edges}.
\end{enumerate}

No método \lstinline{add_edge(u, v, cap, rev_cap)}, os arcos direto e reverso são alocados em pares sequenciais contíguos (índices $2k$ e $2k+1$). Consequentemente, o acesso ao arco residual simétrico opera em tempo estritamente $\mathcal{O}(1)$ via operador bit a bit Ou Exclusivo (\lstinline{id ^ 1ULL}):
\begin{lstlisting}[language=C++, numbers=none, basicstyle=\ttfamily\footnotesize, frame=none, backgroundcolor=\color{white}]
void push_flow(const Size edge_id, const Long flow_amount) {
    edges[edge_id].flow += flow_amount;
    edges[edge_id ^ 1ULL].flow -= flow_amount;
}
\end{lstlisting}

\paragraph{Hierarquia Polimórfica e Padrões de Projeto.}
A biblioteca organiza-se em duas famílias baseadas nos padrões \textit{Strategy}, \textit{Factory Method} e \textit{Prototype}~\cite{gamma1994design}:
\begin{itemize}
    \item \textbf{Família \lstinline{FlowNetwork} (Fluxo Máximo):} Interface base abstrata com o método virtual puro \lstinline{compute_max_flow(source, sink)}. Os motores concretos (\lstinline{FordFulkerson}, \lstinline{EdmondsKarp}, \lstinline{Dinic}, \lstinline{PushRelabel} e \lstinline{PushRelabelImproved}) implementam estratégias de busca distintas mantendo total intercambialidade;
    \item \textbf{Família \lstinline{CostNetwork} (Custo Mínimo):} Interface base abstrata que incorpora o custo linear unitário $w$ aos arcos, com o método virtual puro \lstinline{compute_min_cost(source, sink)} implementado por \lstinline{CycleCanceling}, \lstinline{SuccessiveShortest}, \lstinline{SuccessiveShortestDijkstra} e \lstinline{NetworkSimplex}.
\end{itemize}

Os métodos virtuais \lstinline{make(n)} e \lstinline{clone()} permitem a instanciação e a clonagem profunda (\textit{deep copy}) de instâncias polimórficas preservando o estado residual, o que é fundamental em aplicações iterativas como corte mínimo e segmentação de imagens.

\input{apendices/01_flownetwork.tex}
\input{apendices/02_ford_fulkerson.tex}
\input{apendices/03_edmonds_karp.tex}
\input{apendices/04_dinic.tex}
\input{apendices/05_push_relabel.tex}
\input{apendices/06_push_relabel_improved.tex}
% =========================================================================
% IMPLEMENTAÇÃO DA INTERFACE BASE ``COST-NETWORK''
% =========================================================================
\apendice{Implementação da Interface Base ``Cost-Network''}
\label{ap:costnetwork}

\begin{lstlisting}[caption={Implementação em C++ da interface base \lstinline{CostNetwork} para a modelagem abstrata de redes de fluxo com custos.}, label={lst:costnetwork}]
#ifndef COST_NETWORK_HPP
#define COST_NETWORK_HPP

#include <limits>
#include <memory>
#include <vector>

using Long = long long;
using Size = std::size_t;

constexpr Long INF = std::numeric_limits<Long>::max() >> 8;
constexpr Size MAX = std::numeric_limits<Size>::max() >> 8;

class CostNetwork
{
public:
	struct Edge
	{
		Size from, to;
		Long capacity, cost, flow;
	};

	explicit CostNetwork(const Size n) : size(n), adjacency(n) {}
	virtual ~CostNetwork() = default;

	virtual std::unique_ptr<CostNetwork> make(const Size n) const = 0;
	virtual std::unique_ptr<CostNetwork> clone() const = 0;

	virtual void add_edge(
	    const Size from, const Size to, const Long capacity, const Long cost
	)
	{
		adjacency[from].push_back(edges.size());
		edges.push_back({from, to, capacity, cost, 0});
		adjacency[to].push_back(edges.size());
		edges.push_back({to, from, 0, -cost, 0});
	}

	virtual Long compute_min_cost_max_flow(const Size source, const Size sink) = 0;

	[[nodiscard]] Long get_total_flow(const Size source) const
	{
		Long total = 0;
		for (const Size edge_id : adjacency[source])
			total += edges[edge_id].flow;
		return total;
	}

	[[nodiscard]] Size get_size() const
	{
		return size;
	}

	[[nodiscard]] const std::vector<Edge> &get_edges() const
	{
		return edges;
	}

	[[nodiscard]] const std::vector<std::vector<Size>> &get_adjacency() const
	{
		return adjacency;
	}

protected:
	Size size;
	std::vector<Edge> edges;
	std::vector<std::vector<Size>> adjacency;

	[[nodiscard]] Long get_residual_capacity(const Size edge_id) const
	{
		return edges[edge_id].capacity - edges[edge_id].flow;
	}

	void push_flow(const Size edge_id, const Long flow_amount)
	{
		edges[edge_id].flow += flow_amount;
		edges[edge_id ^ 1ULL].flow -= flow_amount;
	}
};

#endif // COST_NETWORK_HPP
\end{lstlisting}

\input{apendices/08_cycle_canceling.tex}
% =========================================================================
% IMPLEMENTAÇÃO DA CLASSE ``SUCCESSIVE-SHORTEST-PATH''
% =========================================================================
\apendice{Implementação da Classe ``Successive-Shortest-Path''}
\label{ap:successive_shortest}

\begin{lstlisting}[caption={Implementação em C++ do algoritmo \textit{Successive Shortest Path} com SPFA.}, label={lst:successive_shortest}]
#ifndef SUCCESSIVE_SHORTEST_HPP
#define SUCCESSIVE_SHORTEST_HPP

#include "CostNetwork.hpp"
#include <algorithm>
#include <queue>

class SuccessiveShortest : public CostNetwork
{
public:
	explicit SuccessiveShortest(const Size n)
	    : CostNetwork(n), distance(n), parent_edge(n), in_queue(n)
	{
	}

	static std::unique_ptr<CostNetwork> create(const Size n)
	{
		return std::make_unique<SuccessiveShortest>(n);
	}

	std::unique_ptr<CostNetwork> make(const Size n) const override
	{
		return std::make_unique<SuccessiveShortest>(n);
	}

	std::unique_ptr<CostNetwork> clone() const override
	{
		return std::make_unique<SuccessiveShortest>(*this);
	}

	Long compute_min_cost_max_flow(const Size source, const Size sink) override
	{
		Long total_cost = 0;

		while (spfa(source, sink))
		{
			Long bottleneck = INF;
			for (Size node = sink; node != source;)
			{
				const Size edge_id = parent_edge[node];
				bottleneck = std::min(bottleneck, get_residual_capacity(edge_id));
				node = edges[edge_id].from;
			}

			for (Size node = sink; node != source;)
			{
				const Size edge_id = parent_edge[node];
				push_flow(edge_id, bottleneck);
				node = edges[edge_id].from;
			}

			total_cost += bottleneck * distance[sink];
		}

		return total_cost;
	}

private:
	std::vector<Long> distance;
	std::vector<Size> parent_edge;
	std::vector<bool> in_queue;

	bool spfa(const Size source, const Size sink)
	{
		std::fill(distance.begin(), distance.end(), INF);
		std::fill(parent_edge.begin(), parent_edge.end(), MAX);
		std::fill(in_queue.begin(), in_queue.end(), false);

		distance[source] = 0;
		in_queue[source] = true;
		std::queue<Size> queue;
		queue.push(source);

		while (!queue.empty())
		{
			const Size current = queue.front();
			queue.pop();
			in_queue[current] = false;

			for (const Size edge_id : adjacency[current])
			{
				const Size next = edges[edge_id].to;
				const Long new_distance = distance[current] + edges[edge_id].cost;

				if (get_residual_capacity(edge_id) <= 0 ||
				    new_distance >= distance[next])
					continue;

				distance[next] = new_distance;
				parent_edge[next] = edge_id;

				if (!in_queue[next])
				{
					in_queue[next] = true;
					queue.push(next);
				}
			}
		}

		return distance[sink] != INF;
	}
};

#endif // SUCCESSIVE_SHORTEST_HPP
\end{lstlisting}

\input{apendices/10_ssp_dijkstra.tex}
\input{apendices/11_network_simplex.tex}
% =========================================================================
% ACESSO AO CÓDIGO-FONTE E REPOSITÓRIO GITHUB
% =========================================================================
\newpage
\apendice{Acesso ao Código-Fonte e Repositório GitHub}
\label{ap:github}

Para garantir a reprodutibilidade dos experimentos e disponibilizar o código-fonte desenvolvido, todo o material está organizado em repositório público no GitHub.

O repositório \textbf{Networks Flow} está estruturado nos seguintes diretórios principais:
\begin{itemize}[leftmargin=1.5em, itemsep=0.2em]
\raggedright
    \item \textbf{Fluxo Máximo} (\path{Implementações/FlowNetwork/}): Módulos base e algoritmos (\lstinline{FordFulkerson}, \lstinline{EdmondsKarp}, \lstinline{Dinic}, \lstinline{PushRelabel} e \lstinline{PushRelabelImproved});
    \item \textbf{Fluxo de Custo Mínimo} (\path{Implementações/CostNetwork/}): Módulos base e algoritmos (\lstinline{CycleCanceling}, \lstinline{SuccessiveShortestPath} e \lstinline{NetworkSimplex});
    \item \textbf{Problemas} (\path{Implementações/Problemas/}): Resoluções para problemas combinatórios clássicos (como \textit{Minimum Vertex Cover}, emparelhamentos e caminhos disjuntos);
    \item \textbf{Aplicações Práticas} (\path{Aplicações/Segmentação de Imagens/}): Segmentação de imagens via corte mínimo em grafos (\textit{Graph Cuts});
    \item \textbf{Experimentos} (\path{Experimentos/}): Scripts de automação de testes, leitura de instâncias DIMACS e avaliação experimental.
\end{itemize}

O código-fonte completo pode ser consultado, clonado e executado através do endereço:

\vspace{0.4cm}
\begin{center}
    \fbox{\url{https://github.com/GabrielFrigo4/networks-flow}}
\end{center}
\vspace{0.4cm}

O repositório contém o histórico de desenvolvimento e as instruções para compilação e reprodução dos experimentos.

